% Part: normal-modal-logic
% Chapter: syntax-and-semantics
% Section: language-modal-logic

\documentclass[../../../include/open-logic-section]{subfiles}

\begin{document}

\olfileid[hi]{nml}{syn}{lan}

\olsection{मूल मोडल तर्कशास्त्र की भाषा}

\begin{defn}
मोडल तर्कशास्त्र की मूल भाषा में निम्नलिखित होते हैं:
\begin{enumerate}
  \tagitem{prvFalse}{!!{falsity} का प्रतिज्ञप्तिक अचर~$\lfalse$।}{}
  \tagitem{prvTrue}{!!{truth} का प्रतिज्ञप्तिक अचर~$\ltrue$।}{}
  \item !!{propositional variable}ों का !!{denumerable} समुच्चय: $\Obj
    p_0$, $\Obj p_1$, $\Obj p_2$, \dots
  \item प्रतिज्ञप्तिक संयोजक: \startycommalist
  \iftag{prvNot}{\ycomma $\lnot$ (निषेध)}{}%
  \iftag{prvAnd}{\ycomma $\land$ (संयोजन)}{}%
  \iftag{prvOr}{\ycomma $\lor$ (वियोजन)}{}%
  \iftag{prvIf}{\ycomma $\lif$ (!!{conditional})}{}%
  \iftag{prvIff}{\ycomma $\liff$ (!!{biconditional})}{}।
  \tagitem{prvBox}{मोडल संकारक $\Box$।}{}
  \tagitem{prvDiamond}{मोडल संकारक $\Diamond$।}{}
\end{enumerate}
\end{defn}

\begin{defn}
मूल मोडल भाषा की \emph{!!{formula}ें} आगमनात्मक रूप से इस प्रकार परिभाषित
होती हैं:
\begin{enumerate}
\tagitem{prvFalse}{$\lfalse$ एक परमाण्विक !!{formula} है।}{}

\tagitem{prvTrue}{$\ltrue$ एक परमाण्विक !!{formula} है।}{}

\item प्रत्येक !!{propositional variable} $\Obj p_i$ एक (परमाण्विक) !!{formula} है।

\tagitem{prvNot}{यदि $!A$ एक !!{formula} है, तो $\lnot !A$ एक
  !!{formula} है।}{}

\tagitem{prvAnd}{यदि $!A$ और $!B$ !!{formula} हैं, तो $(!A \land
  !B)$ एक !!{formula} है।}{}

\tagitem{prvOr}{यदि $!A$ और $!B$ !!{formula} हैं, तो $(!A \lor !B)$
  एक !!{formula} है।}{}

\tagitem{prvIf}{यदि $!A$ और $!B$ !!{formula} हैं, तो $(!A \lif !B)$
  एक !!{formula} है।}{}

\tagitem{prvIff}{यदि $!A$ और $!B$ !!{formula} हैं, तो $(!A \liff !B)$
  एक !!{formula} है।}{}

\tagitem{prvBox}{यदि $!A$ एक !!{formula} है, तो $\Box !A$ एक
  !!{formula} है।}{}

\tagitem{prvDiamond}{यदि $!A$ एक !!{formula} है, तो $\Diamond !A$ एक
  !!{formula} है।}{}

\tagitem{limitClause}{इसके अतिरिक्त कुछ भी !!{formula} नहीं है।}{}
\end{enumerate}
\end{defn}

\begin{tagblock}{defNot,defOr,defAnd,defIf,defIff,defTrue,defFalse,defBox,defDiamond}
\begin{defn}
परिभाषित संकारकों से निर्मित सूत्रों को इस प्रकार समझना है:

\begin{tagenumerate}{defTrue,defFalse,defNot,defOr,defAnd,defIf,defIff,defBox,defDiamond}

\tagitem{defTrue}{$\ltrue$ इसका संक्षिप्त रूप है:
  \iftag{prvFalse}{$\lnot\lfalse$}{$(!A \lor \lnot !A)$,
  किसी नियत परमाण्विक !!{formula}~$!A$ के लिए}.}{}

\tagitem{defFalse}{$\lfalse$ इसका संक्षिप्त रूप है:
  \iftag{prvTrue}{$\lnot\ltrue$}{$(!A \land \lnot !A)$,
  किसी नियत परमाण्विक !!{formula}~$!A$ के लिए}.}{}

\tagitem{defNot}{$\lnot !A$, $!A \lif \lfalse$ का संक्षिप्त रूप है।}{}

\tagitem{defOr}{$!A \lor !B$ इसका संक्षिप्त रूप है:
  \iftag{prvAnd}{$\lnot(\lnot !A \land \lnot !B)$}{$\lnot !A \lif
    !B$}.}{}

\tagitem{defAnd}{$!A \land !B$ इसका संक्षिप्त रूप है:
  \iftag{prvOr}{$\lnot(\lnot !A \lor \lnot !B)$}{$\lnot (!A \lif
    \lnot !B)$}.}{}

\tagitem{defIf}{$!A \lif !B$ इसका संक्षिप्त रूप है:
  \iftag{prvOr}{$\lnot !A \lor !B)$}{$\lnot (!A \land \lnot !B)$}.}{}

\tagitem{defIff}{$!A \liff !B$, $(!A \lif !B) \land (!B
  \lif !A)$ का संक्षिप्त रूप है।}{}

\tagitem{defBox}{$\Box !A$, $\lnot\Diamond\lnot !A$ का संक्षिप्त रूप है।}{}

\tagitem{defDiamond}{$\Diamond !A$, $\lnot\Box\lnot !A$ का संक्षिप्त रूप है।}{}
\end{tagenumerate}
\end{defn}
\end{tagblock}

यदि किसी !!{formula}~$!A$ में $\Box$ या $\Diamond$ नहीं आता, तो हम उसे
\emph{मोडल-मुक्त} कहते हैं।

\end{document}
